\subsection{Maintenance Effort}
\label{subsec:wartungsaufwand}

The single-repository strategy consolidates the core, features, profiles,
generator, tests, and documentation. This avoids parallel baseline changes and
benefits future projects. Repositories that have already been generated must,
however, incorporate fixes independently.
The source of truth thus reduces the number of starting points that require
maintenance, but not the effort needed to migrate existing products.

Central maintenance encompasses several languages, dependency managers,
toolchains, platforms, and release cycles. This diversity follows from the
web, backend, desktop, and deployment objectives, but increases the expertise
required, the breadth of testing, and dependency on the environment.
Native builds and external services are more maintenance-intensive than a
purely web-based, technologically homogeneous starter.

Five profiles, six features, and one selectable capability currently limit the
combination space; acceptance covers five base paths and three permitted
PostgreSQL extensions \parencite{kleiveist2026acceptance}. Additional
capabilities expand both the dependencies and the test matrix. A combinatorial
explosion has not been demonstrated, but increasing maintenance effort is
plausible.
With such extensions, the catalog, resolver, configuration contract, test
matrix, and documentation must remain mutually consistent.

Central maintenance is therefore more demanding than maintenance of a small,
homogeneous starter, but can consolidate repeated maintenance of the baseline.
Its economic viability depends on frequency of use, similarity between
projects, and maintenance time; cost data are unavailable.
\beginnerinput{workspace/statement/700_bewertung/740}
